\begin{lemma}
Let $\Omega$ and $U$ be open subsets of $\mathbb{C}$ and let
$\chi,H\colon\mathbb{C}\to\mathbb{C}$ be such that
\begin{enumerate}
\item[(a)] $H$ is $C^\infty$ on $\Omega$ and $\chi$ is $C^\infty$ on $U$;
\item[(b)] $H(\xi)\in U$ and $\chi(H(\xi))=\xi$ for every $\xi\in\Omega$;
\item[(c)] $J(\xi)\neq0$ for every $\xi\in\Omega$, where
$J=|H_\xi|^2-|H_{\bar\xi}|^2$ is the Jacobian of \noderef{Jac}.
\end{enumerate}
Let $m\ge1$ be an integer and let $w$ be a word of length $m$ in the two real
partial derivatives $\partial/\partial x$ and $\partial/\partial y$, so that
$D^{w}$ is a partial derivative operator of order $m$ in the sense of
\noderef{PartialWord}.  Then there exists
\[
  P=P_{w}\in\mathbb{P}(4m-3,\,3m-2)
\]
relative to $H$, in the sense of \noderef{PClass}, such that
\[
  D^{w}\chi(z)=J(\xi)^{1-2m}P_{w}(\xi)
  \qquad\text{for every }\xi\in\Omega,\ \text{where }z=H(\xi).
\]
\end{lemma}
\begin{proof}
Throughout, $\Omega$, $U$, $\chi$, $H$ are fixed as in the hypotheses, and all
class memberships are relative to $H$.  Write
\[
  H_\xi=\frac{\partial H}{\partial\xi},\qquad
  H_{\bar\xi}=\frac{\partial H}{\partial\bar\xi}
\]
for the iterated derivatives of \noderef{WirtingerIter} with indices $(1,0)$
and $(0,1)$, and let $\mathbf{J}$ denote the Jacobian $J$ of \noderef{Jac}
regarded as a complex-valued function, $\mathbf{J}(\xi)=J(\xi)\in\mathbb{R}
\subset\mathbb{C}$.  By hypothesis (c), $\mathbf{J}(\xi)\neq0$ for
$\xi\in\Omega$, so integer powers of $\mathbf{J}(\xi)$ are defined there.

\medskip
\noindent\textbf{Step 0.  Five preliminary facts.}

\emph{(0a) $\mathbb{P}(s,k)$ is a $\mathbb{C}$-linear space.}  If
$P,Q\in\mathbb{P}(s,k)$ and $c\in\mathbb{C}$ then $P+Q\in\mathbb{P}(s,k)$ and
$cP\in\mathbb{P}(s,k)$.  This is immediate from \noderef{PClass}: if
$(c_\alpha,M_\alpha)_{\alpha\le r}$ presents $P$ and $(b_\beta,N_\beta)_{\beta\le q}$
presents $Q$, with all lists of type $(s,k)$, then the concatenated family
presents $P+Q$ and the family $(c\,c_\alpha,M_\alpha)_{\alpha\le r}$ presents
$cP$, and in both cases every list in the new family is one of the old ones and
so still has type $(s,k)$.

\emph{(0b) The four first-order factors lie in $\mathbb{P}(1,1)$.}  Consider
the derivative symbols of \noderef{DerivSymbol}
\[
  e=(1,0,\text{false}),\quad e^{c}=(1,0,\text{true}),\quad
  e'=(0,1,\text{false}),\quad e'^{c}=(0,1,\text{true}),
\]
each of order $1$.  By \noderef{SymbolValue} their values on $H$ are
$H_\xi$, $\overline{H_\xi}$, $H_{\bar\xi}$, $\overline{H_{\bar\xi}}$
respectively.  Each of the one-entry lists $(e)$, $(e^{c})$, $(e')$,
$(e'^{c})$ has length $1$ and total order $1$, i.e.\ type $(1,1)$, and by
\noderef{MonomialValue} its monomial is the corresponding symbol value.  Taking
the single coefficient $1$ in \noderef{PClass},
\[
  H_\xi,\ \overline{H_\xi},\ H_{\bar\xi},\ \overline{H_{\bar\xi}}\ \in\ \mathbb{P}(1,1).
\]

\emph{(0c) $\mathbf{J}\in\mathbb{P}(2,2)$, $\mathbf{J}$ is $C^\infty$ on
$\Omega$, and both $\partial\mathbf{J}/\partial\xi$ and
$\partial\mathbf{J}/\partial\bar\xi$ agree on $\Omega$ with members of
$\mathbb{P}(3,2)$.}  The membership $\mathbf{J}\in\mathbb{P}(2,2)$ is part (a)
of \noderef{JacInPClass}.  Hence $\mathbf{J}$ is $C^\infty$ on $\Omega$ by
\noderef{PClassSmooth}, using hypothesis (a) and that $\Omega$ is open.  Part
(b) of \noderef{JacInPClass} provides a function $R\in\mathbb{P}(3,2)$, namely
\[
  R=H_{\xi\xi}\overline{H_\xi}+\overline{H_{\xi\bar\xi}}H_\xi
     -H_{\xi\bar\xi}\overline{H_{\bar\xi}}-\overline{H_{\bar\xi\bar\xi}}H_{\bar\xi},
\]
and part (c) states
\begin{equation}\tag{1}
  \frac{\partial\mathbf{J}}{\partial\xi}(\xi)=R(\xi)\qquad\text{for every }\xi\in\Omega .
\end{equation}
For the other operator we use that $J$ is real-valued.  By \noderef{Jac},
$J(\xi)=\|H_\xi(\xi)\|^2-\|H_{\bar\xi}(\xi)\|^2$ is a difference of squared
norms, hence a real number, so the complex number $\mathbf{J}(\xi)$ satisfies
$\overline{\mathbf{J}(\xi)}=\mathbf{J}(\xi)$ for every $\xi$; that is, the
function $\overline{\mathbf{J}}$ equals $\mathbf{J}$.  Since $\mathbf{J}$ is
$C^\infty$ on $\Omega$ it is real-differentiable at each $\xi\in\Omega$, so
\noderef{WirtingerConjRule} applies to $\mathbf{J}$ at $\xi$ and gives
\[
  \frac{\partial\overline{\mathbf{J}}}{\partial\bar\xi}(\xi)
  =\overline{\frac{\partial\mathbf{J}}{\partial\xi}(\xi)} .
\]
The left-hand side is $\partial\mathbf{J}/\partial\bar\xi(\xi)$, because
$\overline{\mathbf{J}}$ and $\mathbf{J}$ are the same function.  Combining with
(1),
\begin{equation}\tag{2}
  \frac{\partial\mathbf{J}}{\partial\bar\xi}(\xi)=\overline{R(\xi)}
  \qquad\text{for every }\xi\in\Omega ,
\end{equation}
and $\overline{R}\in\mathbb{P}(3,2)$ by \noderef{PClassConjClosed}.

\emph{(0d) For every word $v$, $D^{v}\chi$ is $C^\infty$ on $U$.}  By
hypothesis (a), $\chi$ is $C^\infty$ on the open set $U$.  By
\noderef{PartialWord}, $D^{v}\chi$ is obtained from $\chi$ by applying $|v|$ of
the two real first-order operators $\partial/\partial x$,
$\partial/\partial y$, and each of these carries $C^\infty(U)$ into
$C^\infty(U)$, because the first partial derivatives of a $C^\infty$ function
on an open set are again $C^\infty$ there.  Induction on $|v|$ gives the claim.
In particular $D^{v}\chi$ is real-differentiable at every point of $U$, hence
at $H(\xi)$ for every $\xi\in\Omega$ by hypothesis (b).

\emph{(0e) The two real partials in terms of the Wirtinger operators.}  Let $F$
be real-differentiable at a point $\zeta$ and write $F_x(\zeta)=DF(\zeta)[1]$,
$F_y(\zeta)=DF(\zeta)[i]$ for its two real first partials.  By \noderef{DZ} and
\noderef{DZbar},
\[
  \frac{\partial F}{\partial z}(\zeta)+\frac{\partial F}{\partial\bar z}(\zeta)
  =\tfrac12\bigl(F_x-iF_y\bigr)(\zeta)+\tfrac12\bigl(F_x+iF_y\bigr)(\zeta)
  =F_x(\zeta),
\]
\[
  \frac{\partial F}{\partial z}(\zeta)-\frac{\partial F}{\partial\bar z}(\zeta)
  =-i\,F_y(\zeta),
  \qquad\text{hence}\qquad
  i\Bigl(\frac{\partial F}{\partial z}(\zeta)-\frac{\partial F}{\partial\bar z}(\zeta)\Bigr)
  =i\bigl(-i\,F_y(\zeta)\bigr)=F_y(\zeta).
\]
By \noderef{PartialWord}, for a word $v$ and a direction $j\in\{0,1\}$ the word
$(j)*v$ obtained by prepending $j$ satisfies
$D^{(j)*v}F=\partial_j\bigl(D^{v}F\bigr)$ with $\partial_0=\partial/\partial x$
and $\partial_1=\partial/\partial y$.  Applying the two displays to
$F=D^{v}\chi$, which is real-differentiable on $U$ by (0d), gives at every
point of $U$
\begin{equation}\tag{3}
  D^{(0)*v}\chi=\bigl(D^{v}\chi\bigr)_z+\bigl(D^{v}\chi\bigr)_{\bar z},
  \qquad
  D^{(1)*v}\chi=i\Bigl(\bigl(D^{v}\chi\bigr)_z-\bigl(D^{v}\chi\bigr)_{\bar z}\Bigr).
\end{equation}

\medskip
\noindent\textbf{Step 1.  The base case $m=1$.}

Let $w=(j)$ be a word of length $1$ and let $\xi\in\Omega$.  We apply
\noderef{ChiFirstDerivs} with its open set taken to be $\Omega$: it contains
$\xi$; the inverse relation $\chi(H(\zeta))=\zeta$ holds for every
$\zeta\in\Omega$ by hypothesis (b); $H$ is real-differentiable at $\xi$ by
hypothesis (a) and openness of $\Omega$; $\chi$ is real-differentiable at
$H(\xi)$ by hypotheses (a), (b) and openness of $U$; and $J(\xi)\neq0$ by
hypothesis (c).  It yields
\[
  \frac{\partial\chi}{\partial z}\bigl(H(\xi)\bigr)=\frac{\overline{H_\xi(\xi)}}{\mathbf{J}(\xi)},
  \qquad
  \frac{\partial\chi}{\partial\bar z}\bigl(H(\xi)\bigr)=-\frac{H_{\bar\xi}(\xi)}{\mathbf{J}(\xi)} .
\]
Taking $v$ to be the empty word in (3), for which $D^{v}\chi=\chi$ by
\noderef{PartialWord}, and evaluating at the point $H(\xi)\in U$,
\[
  D^{(0)}\chi\bigl(H(\xi)\bigr)
  =\mathbf{J}(\xi)^{-1}\Bigl(\overline{H_\xi(\xi)}-H_{\bar\xi}(\xi)\Bigr),
  \qquad
  D^{(1)}\chi\bigl(H(\xi)\bigr)
  =\mathbf{J}(\xi)^{-1}\Bigl(i\,\overline{H_\xi(\xi)}+i\,H_{\bar\xi}(\xi)\Bigr),
\]
the second because
$i\bigl(\overline{H_\xi}/\mathbf{J}-(-H_{\bar\xi}/\mathbf{J})\bigr)
 =\mathbf{J}^{-1}\bigl(i\overline{H_\xi}+iH_{\bar\xi}\bigr)$.
Put
\[
  P_{(0)}=\overline{H_\xi}-H_{\bar\xi},
  \qquad
  P_{(1)}=i\,\overline{H_\xi}+i\,H_{\bar\xi} .
\]
By (0b) and (0a) both belong to $\mathbb{P}(1,1)$, and
$\mathbb{P}(1,1)=\mathbb{P}(4\cdot1-3,\,3\cdot1-2)$.  Since
$\mathbf{J}(\xi)^{-1}=\mathbf{J}(\xi)^{1-2\cdot1}$, the conclusion holds for
$m=1$ with $P=P_{w}$.

\medskip
\noindent\textbf{Step 2.  The inductive step.}

Let $m\ge1$ and assume the conclusion for every word of length $m$: for each
such word $v$ there is $P_{v}\in\mathbb{P}(4m-3,3m-2)$ with
\begin{equation}\tag{IH}
  D^{v}\chi\bigl(H(\xi)\bigr)=\mathbf{J}(\xi)^{1-2m}P_{v}(\xi)
  \qquad\text{for every }\xi\in\Omega .
\end{equation}
Let $w'$ be a word of length $m+1$ and write $w'=(j)*v$ with $|v|=m$.

We claim it suffices to produce $Q_1,Q_2\in\mathbb{P}(4m+1,3m+1)$ with
\begin{equation}\tag{4}
  \bigl(D^{v}\chi\bigr)_z\bigl(H(\xi)\bigr)=\mathbf{J}(\xi)^{-1-2m}Q_1(\xi),
  \qquad
  \bigl(D^{v}\chi\bigr)_{\bar z}\bigl(H(\xi)\bigr)=\mathbf{J}(\xi)^{-1-2m}Q_2(\xi)
\end{equation}
for every $\xi\in\Omega$.  Indeed, granting (4), equation (3) evaluated at
$H(\xi)$ gives
\[
  D^{w'}\chi\bigl(H(\xi)\bigr)=
  \begin{cases}
    \mathbf{J}(\xi)^{-1-2m}\bigl(Q_1(\xi)+Q_2(\xi)\bigr)&\text{if }j=0,\\
    \mathbf{J}(\xi)^{-1-2m}\bigl(iQ_1(\xi)-iQ_2(\xi)\bigr)&\text{if }j=1,
  \end{cases}
\]
and in either case the second factor is a $\mathbb{C}$-linear combination of
$Q_1$ and $Q_2$, hence lies in $\mathbb{P}(4m+1,3m+1)$ by (0a).  Since
$m\ge1$,
\[
  -1-2m=1-2(m+1),\qquad 4m+1=4(m+1)-3,\qquad 3m+1=3(m+1)-2,
\]
so this is exactly the conclusion for the word $w'$ of length $m+1$.

\emph{Step 2a: the chain-rule system.}  Let $\xi\in\Omega$ be arbitrary; it is
fixed only for Steps 2a--2c, and the functions $Q_1,Q_2$ constructed in Step 2d
will not depend on it.  Abbreviate
\[
  a=\bigl(D^{v}\chi\bigr)_z\bigl(H(\xi)\bigr),\quad
  b=\bigl(D^{v}\chi\bigr)_{\bar z}\bigl(H(\xi)\bigr),\quad
  A=H_\xi(\xi),\quad B=H_{\bar\xi}(\xi).
\]
By (0d), $D^{v}\chi$ is real-differentiable at $H(\xi)$, and by hypothesis (a)
with $\Omega$ open, $H$ is real-differentiable at $\xi$.  Hence
\noderef{WirtingerChainRule} applies with $F=D^{v}\chi$ and $G=H$ and gives
\begin{equation}\tag{5}
  \frac{\partial\bigl(D^{v}\chi\circ H\bigr)}{\partial\xi}(\xi)=a\,A+b\,\overline{B},
  \qquad
  \frac{\partial\bigl(D^{v}\chi\circ H\bigr)}{\partial\bar\xi}(\xi)=a\,B+b\,\overline{A}.
\end{equation}
This is the pair of equations displayed in the proof of Lemma 1 in the paper.

\emph{Step 2b: computing the left-hand sides.}  By (IH) the function
$\zeta\mapsto D^{v}\chi(H(\zeta))$ agrees at every point of $\Omega$ with
$\zeta\mapsto\mathbf{J}(\zeta)^{1-2m}P_{v}(\zeta)$.  Since $\Omega$ is open and
the real Fr\'echet derivative at $\xi$ depends only on the restriction of a
function to a neighbourhood of $\xi$, the two functions have the same
$\partial/\partial\xi$ and $\partial/\partial\bar\xi$ at $\xi$, by \noderef{DZ}
and \noderef{DZbar}.

Both factors are $C^\infty$ on $\Omega$: $P_{v}$ by \noderef{PClassSmooth}
applied to $P_{v}\in\mathbb{P}(4m-3,3m-2)$, and $\zeta\mapsto\mathbf{J}(\zeta)^{1-2m}$
because $\mathbf{J}$ is $C^\infty$ on $\Omega$ by (0c) and non-vanishing there
by hypothesis (c), so that a negative integer power of it is a finite product
of copies of the $C^\infty$ function $1/\mathbf{J}$.  In particular both are
real-differentiable at $\xi$.  By \noderef{WirtingerProductRule} and then
\noderef{WirtingerPowerRule} applied to $g=\mathbf{J}$ with exponent $a=1-2m$
(legitimate by (0c) and hypothesis (c)),
\[
  \frac{\partial\bigl(\mathbf{J}^{1-2m}P_{v}\bigr)}{\partial\xi}(\xi)
  =(1-2m)\,\mathbf{J}(\xi)^{-2m}\frac{\partial\mathbf{J}}{\partial\xi}(\xi)\,P_{v}(\xi)
   +\mathbf{J}(\xi)^{1-2m}\frac{\partial P_{v}}{\partial\xi}(\xi),
\]
using $(1-2m)-1=-2m$.  Since $\mathbf{J}(\xi)\neq0$ we may write
$\mathbf{J}(\xi)^{1-2m}=\mathbf{J}(\xi)^{-2m}\mathbf{J}(\xi)$ and factor out
$\mathbf{J}(\xi)^{-2m}$.  Writing $L_1$ and $L_2$ for the two left-hand sides
of (5), we obtain
\begin{equation}\tag{6}
  L_1=\mathbf{J}(\xi)^{-2m}\Bigl[\mathbf{J}(\xi)\frac{\partial P_{v}}{\partial\xi}(\xi)
      +(1-2m)P_{v}(\xi)\frac{\partial\mathbf{J}}{\partial\xi}(\xi)\Bigr],
\end{equation}
\begin{equation}\tag{7}
  L_2=\mathbf{J}(\xi)^{-2m}\Bigl[\mathbf{J}(\xi)\frac{\partial P_{v}}{\partial\bar\xi}(\xi)
      +(1-2m)P_{v}(\xi)\frac{\partial\mathbf{J}}{\partial\bar\xi}(\xi)\Bigr],
\end{equation}
(7) by the identical computation with $\partial/\partial\bar\xi$ in place of
$\partial/\partial\xi$, using the second identity of
\noderef{WirtingerProductRule} and the second identity of
\noderef{WirtingerPowerRule}.  Equations (5)--(7) are the system
\[
  \left\{
  \begin{aligned}
    a\,H_\xi+b\,\overline{H_{\bar\xi}}&=\mathbf{J}^{-2m}\bigl[\mathbf{J}P_\xi+(1-2m)P\mathbf{J}_\xi\bigr]\\
    a\,H_{\bar\xi}+b\,\overline{H_\xi}&=\mathbf{J}^{-2m}\bigl[\mathbf{J}P_{\bar\xi}+(1-2m)P\mathbf{J}_{\bar\xi}\bigr]
  \end{aligned}
  \right.
\]
of the paper, evaluated at $\xi$, with $P=P_v$.

\emph{Step 2c: solving the system.}  First the determinant.  For
$\omega\in\mathbb{C}$ one has $|\omega|^2=\omega\bar\omega$ as complex numbers,
so by \noderef{Jac}
\begin{equation}\tag{8}
  \mathbf{J}(\xi)=A\overline{A}-B\overline{B}.
\end{equation}
Multiply the first equation of (5) by $\overline{A}$ and the second by
$\overline{B}$ and subtract.  The terms in $b$ cancel, because
$\overline{B}\,\overline{A}=\overline{A}\,\overline{B}$, leaving
\[
  a\bigl(A\overline{A}-B\overline{B}\bigr)=L_1\overline{A}-L_2\overline{B},
  \qquad\text{i.e.}\qquad
  a\,\mathbf{J}(\xi)=L_1\overline{A}-L_2\overline{B}
\]
by (8).  Since $\mathbf{J}(\xi)\neq0$,
\begin{equation}\tag{9}
  a=\mathbf{J}(\xi)^{-1}\bigl(L_1\overline{A}-L_2\overline{B}\bigr).
\end{equation}
Multiply the first equation of (5) by $B$ and the second by $A$ and subtract.
The terms in $a$ cancel, because $AB=BA$, leaving
\[
  b\bigl(\overline{B}B-\overline{A}A\bigr)=L_1B-L_2A,
  \qquad\text{i.e.}\qquad
  -b\,\mathbf{J}(\xi)=L_1B-L_2A ,
\]
again by (8), so
\begin{equation}\tag{10}
  b=\mathbf{J}(\xi)^{-1}\bigl(L_2A-L_1B\bigr).
\end{equation}
Substituting (6) and (7) into (9) and (10), and using
$\mathbf{J}(\xi)^{-1}\mathbf{J}(\xi)^{-2m}=\mathbf{J}(\xi)^{-1-2m}$ (valid
since $\mathbf{J}(\xi)\neq0$),
\begin{equation}\tag{11}
  a=\mathbf{J}(\xi)^{-1-2m}\Bigl[\mathbf{J}P_\xi\overline{A}-\mathbf{J}P_{\bar\xi}\overline{B}
     +(1-2m)P\mathbf{J}_\xi\overline{A}-(1-2m)P\mathbf{J}_{\bar\xi}\overline{B}\Bigr](\xi),
\end{equation}
\begin{equation}\tag{12}
  b=\mathbf{J}(\xi)^{-1-2m}\Bigl[\mathbf{J}P_{\bar\xi}A-\mathbf{J}P_\xi B
     +(1-2m)P\mathbf{J}_{\bar\xi}A-(1-2m)P\mathbf{J}_\xi B\Bigr](\xi),
\end{equation}
where $P=P_v$.  These are exactly the two expressions the paper displays for
$(D^\nu\chi)_z$ and $(D^\nu\chi)_{\bar z}$.  The single extra factor
$\mathbf{J}(\xi)^{-1}$, which turns $\mathbf{J}^{-2m}$ into
$\mathbf{J}^{-1-2m}$, is precisely the division by the determinant of the
system, and the determinant is $\mathbf{J}(\xi)$ by (8).

\emph{Step 2d: every bracket term lies in $\mathbb{P}(4m+1,3m+1)$.}  Collect
the memberships, all relative to $H$.  Each item below produces a function on
$\mathbb{C}$ that does not depend on the point $\xi$ fixed in Step 2a: the
existential quantifiers of \noderef{PClassDerivClosed} and the class members of
\noderef{JacInPClass} stand outside the ``for every $\xi\in\Omega$'' of those
statements.  The list is:
\begin{itemize}
\item $\mathbf{J}\in\mathbb{P}(2,2)$, by (0c);
\item by \noderef{PClassDerivClosed} applied to $P_{v}\in\mathbb{P}(4m-3,3m-2)$
(its hypotheses hold: $\Omega$ is open and $H$ is $C^\infty$ on $\Omega$), there
are
\[
  Q,\,Q'\in\mathbb{P}\bigl((4m-3)+1,\,3m-2\bigr)=\mathbb{P}(4m-2,\,3m-2)
\]
with $\partial P_{v}/\partial\xi=Q$ and $\partial P_{v}/\partial\bar\xi=Q'$ at
every point of $\Omega$;
\item $R,\ \overline{R}\in\mathbb{P}(3,2)$ with
$\partial\mathbf{J}/\partial\xi=R$ and
$\partial\mathbf{J}/\partial\bar\xi=\overline{R}$ on $\Omega$, by (0c);
\item $P_{v}\in\mathbb{P}(4m-3,3m-2)$, by (IH);
\item $H_\xi,\ \overline{H_\xi},\ H_{\bar\xi},\ \overline{H_{\bar\xi}}\in\mathbb{P}(1,1)$,
by (0b).
\end{itemize}
Now \noderef{PClassMulClosed} closes $\mathbb{P}$ under pointwise products,
adding both indices; applying it twice gives the class of a product of three
factors.  The eight products occurring in the brackets of (11) and (12) have
only two shapes, and both land in $\mathbb{P}(4m+1,3m+1)$:
\begin{align*}
  \mathbf{J}\cdot Q\cdot\overline{H_\xi},\quad
  \mathbf{J}\cdot Q'\cdot\overline{H_{\bar\xi}},\quad
  \mathbf{J}\cdot Q'\cdot H_\xi,\quad
  \mathbf{J}\cdot Q\cdot H_{\bar\xi}
  &:\quad
  \bigl(2+(4m-2)+1,\ 2+(3m-2)+1\bigr)=(4m+1,\,3m+1),\\
  P_{v}\cdot R\cdot\overline{H_\xi},\quad
  P_{v}\cdot\overline{R}\cdot\overline{H_{\bar\xi}},\quad
  P_{v}\cdot\overline{R}\cdot H_\xi,\quad
  P_{v}\cdot R\cdot H_{\bar\xi}
  &:\quad
  \bigl((4m-3)+3+1,\ (3m-2)+2+1\bigr)=(4m+1,\,3m+1).
\end{align*}
Both index computations use $m\ge1$, so that $4m-3$ and $3m-2$ are the
ordinary differences: $2+(4m-2)+1=4m+1$ and $2+(3m-2)+1=3m+1$, while
$(4m-3)+3+1=4m+1$ and $(3m-2)+2+1=3m+1$.

Define
\begin{align*}
  Q_1&=\mathbf{J}\cdot Q\cdot\overline{H_\xi}
      -\mathbf{J}\cdot Q'\cdot\overline{H_{\bar\xi}}
      +(1-2m)\,P_{v}\cdot R\cdot\overline{H_\xi}
      -(1-2m)\,P_{v}\cdot\overline{R}\cdot\overline{H_{\bar\xi}},\\
  Q_2&=\mathbf{J}\cdot Q'\cdot H_\xi
      -\mathbf{J}\cdot Q\cdot H_{\bar\xi}
      +(1-2m)\,P_{v}\cdot\overline{R}\cdot H_\xi
      -(1-2m)\,P_{v}\cdot R\cdot H_{\bar\xi},
\end{align*}
as pointwise products and sums of functions on $\mathbb{C}$.  Each of the eight
summands lies in $\mathbb{P}(4m+1,3m+1)$ by the computation just made, and the
scalars $\pm1$ and $\pm(1-2m)$ are complex numbers, so
$Q_1,Q_2\in\mathbb{P}(4m+1,3m+1)$ by (0a).

Finally let $\xi\in\Omega$.  The bracket in (11) and the value $Q_1(\xi)$
differ only in that the bracket contains
$\partial P_{v}/\partial\xi(\xi)$, $\partial P_{v}/\partial\bar\xi(\xi)$,
$\partial\mathbf{J}/\partial\xi(\xi)$, $\partial\mathbf{J}/\partial\bar\xi(\xi)$
where $Q_1(\xi)$ contains $Q(\xi)$, $Q'(\xi)$, $R(\xi)$, $\overline{R(\xi)}$,
and these pairs agree at every point of $\Omega$ by the four identifications
just listed; also $\overline{A}=\overline{H_\xi(\xi)}$ and
$\overline{B}=\overline{H_{\bar\xi}(\xi)}$ by definition of $A$ and $B$.
Hence (11) reads $a=\mathbf{J}(\xi)^{-1-2m}Q_1(\xi)$, and in the same way (12)
reads $b=\mathbf{J}(\xi)^{-1-2m}Q_2(\xi)$.  Since $Q_1$ and $Q_2$ were defined
without reference to $\xi$, and $\xi\in\Omega$ was arbitrary, these two
identities hold for every $\xi\in\Omega$.  This is (4), and the inductive step
is complete.

\medskip
\noindent\textbf{Step 3.  Conclusion.}  By Step 1 the conclusion holds for
every word of length $1$, and by Step 2 its validity for all words of length
$m\ge1$ implies its validity for all words of length $m+1$.  By induction it
holds for every $m\ge1$ and every word $w$ of length $m$, which is the
assertion of the lemma.
\end{proof}
